\section*{第 \ref{ch:g10:functions} 章——函数}

\begin{solution}{exo:g10:functions:1}
$f(0) = 1$; \quad $f(2) = 8 - 6 + 1 = 3$; \quad
$f(-1) = 2 + 3 + 1 = 6$; \quad
$f\!\left(\tfrac12\right) = 2 \times \tfrac14 - \tfrac32 + 1 = \tfrac12 -
\tfrac32 + 1 = 0$.
\end{solution}

\begin{solution}{exo:g10:functions:2}
$f$：无限制条件，\omterm{def:g10:functions:function}{定义域}为 $\R$。

$g$：要求 $x + 4 \neq 0$，\omterm{def:g10:functions:function}{定义域}为除 $-4$ 外的全体\omterm{def:g10:numbers:sets}{实数}。

$h$：要求 $2x - 6 \geq 0$，即 $x \geq 3$，\omterm{def:g10:functions:function}{定义域}为
$\intco{3}{+\infty}$。

$k$：$x^2 + 1 \geq 1 > 0$，永不为零，\omterm{def:g10:functions:function}{定义域}为 $\R$。
\end{solution}

\begin{solution}{exo:g10:functions:3}
\emph{$0$ 的\omterm{def:g10:functions:function}{原像}：}$x^2 - 2x = x(x-2) = 0$，故为 $0$ 和 $2$。

\emph{$3$ 的\omterm{def:g10:functions:function}{原像}：}$x^2 - 2x = 3$，即 $x^2 - 2x - 3 = 0$，亦即
$(x-3)(x+1) = 0$（\omterm{def:g10:algebra:expand}{展开}可检验），故为 $3$ 和 $-1$。

\emph{$-1$ 的\omterm{def:g10:functions:function}{原像}：}$x^2 - 2x + 1 = (x - 1)^2 = 0$，故唯一的
\omterm{def:g10:functions:function}{原像}为 $1$。
\end{solution}

\begin{solution}{exo:g10:functions:4}
$f(-1) = 1 + 2 + 1 = 4$，故点 $A(-1, 4)$ 在\omterm{def:g10:functions:graph}{图像}上。
$f(2) = 4 - 4 + 1 = 1$，故点 $B(2, 1)$ 也在\omterm{def:g10:functions:graph}{图像}上。
\end{solution}

\begin{solution}{exo:g10:functions:5}
\emph{1.} 表中最大的值为 $4 = g(3)$（\omterm{def:g10:functions:extrema}{最大值}），
最小的值为 $-2 = g(0)$（\omterm{def:g10:functions:extrema}{最小值}）。

\emph{2.} 在 $\intcc{-3}{0}$ 上\omterm{def:g10:functions:function}{函数}\omterm{def:g10:functions:variations}{递减}，且
$-1 < -0.5$，故 $g(-1) \geq g(-0.5)$。

\emph{3.} 在每个单调\omterm{def:g10:numbers:interval}{区间}上，$g$ 每个值至多取一次，
故 $g(x) = 0$ 在每个\omterm{def:g10:numbers:interval}{区间}上至多有一个解，因此总共至多 $3$
个解。（此处 $0$ 介于 $-2$ 与 $1$ 之间，介于 $-2$ 与
$4$ 之间，又在 $0$ 与 $4$ 之间，故恰有三个。）
\end{solution}

\begin{solution}{exo:g10:functions:6}
设 $u < v$，则
\[
f(v) - f(u) = (-2v + 7) - (-2u + 7) = -2(v - u) < 0,
\]
因 $v - u > 0$，于是 $f(v) < f(u)$，即 $f$ 在 $\R$ 上严格\omterm{def:g10:functions:variations}{递减}。
\end{solution}

\begin{solution}{exo:g10:functions:7}
配方：$(x+3)^2 = x^2 + 6x + 9$，于是
\[
f(x) = x^2 + 6x + 11 = (x + 3)^2 + 2 .
\]
对每个 $x$，$(x+3)^2 \geq 0$，故 $f(x) \geq 2$；且
$f(-3) = 0 + 2 = 2$。$f$ 的\omterm{def:g10:functions:extrema}{最小值}为 $2$，在 $x = -3$ 处取得。
\end{solution}

\begin{solution}{exo:g10:functions:8}
\emph{1.} 长 $+$ 宽 $= 20$，故长为 $20 - x$，则
\[
A(x) = x(20 - x).
\]
长与宽均须为正：$0 < x < 20$。

\emph{2.} $A(4) = 64$，$A(8) = 96$，$A(10) = 100$，$A(12) = 96$，
$A(16) = 64$。面积值在 $x = 10$ 之前\omterm{def:g10:functions:variations}{递增}，此后对称\omterm{def:g10:functions:variations}{递减}：
面积似乎在 $x = 10$ 时最大，即当花园为\emph{正方形}时最大。
（此猜想由\cref{ch:g11:quad}证明。）
\end{solution}

\begin{solution}{exo:g10:functions:9}
\emph{1.} 对每个 $x$：$x^2 + 1 \geq 1 > 0$，故 $f(x) = \frac{1}{x^2+1}$
为正，且由 $x^2 + 1 \geq 1$，两边取倒数（两边均为正）
得 $f(x) \leq 1$。

\emph{2.} \omterm{def:g10:functions:extrema}{最大值}：$f(0) = 1$，$f(x) \leq 1$ 对所有 $x$ 成立，故 $1$ 为
\omterm{def:g10:functions:extrema}{最大值}，在 $0$ 处取得。\omterm{def:g10:functions:extrema}{最小值}：$f(x) > 0$ 恒成立，但不存在 $x$
使其取值为 $0$（非零数之商仍非零），且 $f$ 的值可以
任意接近 $0$（当 $x$ 很大时），故 $f$ 没有\omterm{def:g10:functions:extrema}{最小值}。
\end{solution}

\begin{solution}{exo:g10:functions:10}
\emph{1.} $f(v) - f(u) = v^2 - u^2 = (v-u)(v+u)$。若 $0 \leq u < v$，则
$v - u > 0$，且 $v + u > 0$（非负数与正数之和），
故 $f(v) - f(u) > 0$，即 $f$ 在 $\intco{0}{+\infty}$ 上严格\omterm{def:g10:functions:variations}{递增}。

\emph{2.} 若 $u < v \leq 0$，则 $v - u > 0$ 仍成立，但此时
$v + u < 0$（非正数与负数之和），故
$f(v) - f(u) < 0$，即 $f$ 在 $\intoc{-\infty}{0}$ 上严格\omterm{def:g10:functions:variations}{递减}。
\end{solution}

\begin{solution}{pb:g10:functions:1}
\textbf{1.} \omterm{def:g10:functions:function}{定义域}：$x \neq 0$（因需除以 $x$）。$f(1) =
\frac12(1 + 2) = \frac32$; $f(2) = \frac12(2 + 1) = \frac32$;
$f\!\left(\frac32\right) = \frac12\left(\frac32 +
\frac43\right) = \frac{17}{12}$.

\textbf{2.} 由 $\frac12\left(x + \frac2x\right) = \frac32$ 可得
$x^2 + 2 = 3x$，即 $x^2 - 3x + 2 = (x - 1)(x - 2) = 0$：
$\frac32$ 的\omterm{def:g10:functions:function}{原像}为 $1$ 和 $2$——正如第 1 题
所预告的。

\textbf{3.} 由 $f(x) = x$ 得 $x + \frac2x = 2x$，于是
$\frac2x = x$，$x^2 = 2$：\omterm{pb:g10:functions:1}{不动点}为 $\sqrt2$ 和 $-\sqrt2$。
正的那个是单位正方形的对角线，
由初中分册的无理性周末问题可知，它是\omterm{ex:g10:numbers:classify}{无理数}。

\textbf{4.} $1 \to \frac32 \to \frac{17}{12} \to
\frac{577}{408}$：海伦给出的 $\sqrt2$ 的\omterm{def:g10:numbers:approx}{近似值}。他的
做法“将猜测值与 $2/$ 猜测值取平均”恰是
\emph{反复迭代\omterm{def:g10:functions:function}{函数} $f$}——而这一过程所追求的
数，正是 $f$ 的\omterm{pb:g10:functions:1}{不动点}。

\textbf{5.} \omterm{pb:g10:functions:1}{不动点}即\omterm{def:g10:functions:graph}{图像}与
对角线 $y = x$ 的交点：反复将输出值代回，点便
沿\omterm{def:g10:functions:graph}{图像}走向该交点。在
初中分册的“两支温度计”周末问题中，读数 $-40^\circ$ 恰是
换算\omterm{def:g10:functions:function}{函数}“与对角线相交”的同一思想。

\textbf{6.} $f(v) - f(u) = \frac{v - u}{2} + \frac1v - \frac1u
= \frac{v - u}{2} + \frac{u - v}{uv}
= (v - u)\left(\frac12 - \frac{1}{uv}\right)
= \frac{(v - u)(uv - 2)}{2uv}$。当 $0 < u < v \leq \sqrt2$ 时，
$uv < 2$，括号为负，$f(v) < f(u)$，故\omterm{def:g10:functions:variations}{递减}。
当 $\sqrt2 \leq u < v$ 时，$uv > 2$，\omterm{def:g10:functions:variations}{递增}。转折点为
$\sqrt2$。

\textbf{7.} 在 $\sqrt2$ 之前\omterm{def:g10:functions:variations}{递减}，之后\omterm{def:g10:functions:variations}{递增}：$f$
在 $\sqrt2$ 处取得\omterm{def:g10:functions:extrema}{最小值}，其值为
$f(\sqrt2) = \frac12\left(\sqrt2 + \frac{2}{\sqrt2}\right)
= \sqrt2$。由此可知，机器的每个输出（输入为正时）
都不小于 $\sqrt2$——经一次迭代后，
猜测值落在 $\intco{\sqrt2}{+\infty}$ 中。

\textbf{8.} 若 $x > \sqrt2$，则 $x^2 > 2$，于是
$\frac2x < x$，且 $f(x)$ 是 $x$ 与 $\frac2x$ 的平均数，
小于 $x$；由第 7 题可知 $f(x) \geq \sqrt2$，等号
仅在 $x = \sqrt2$ 时成立。因此
$\sqrt2 < f(x) < x$：每次迭代都严格改进而不
过头——第 4 题中的数列沿此下降，趋近
$\sqrt2$。

\textbf{9.} 在 $\intoo{0}{+\infty}$ 上：\omterm{def:g10:functions:function}{函数}先从
$+\infty$（接近 $0$ 时）\omterm{def:g10:functions:variations}{递减}到\omterm{def:g10:functions:extrema}{最小值} $\sqrt2$，
在 $x = \sqrt2$ 处取得，随后无界\omterm{def:g10:functions:variations}{递增}。

\textbf{10.} 由 $g(x) = x$ 得 $x^2 = x$，$x(x - 1) = 0$：\omterm{pb:g10:functions:1}{不动点}为
$0$ 和 $1$。从 $0.9$ 出发：$0.81$，$0.656$，$0.430$，
$0.185$——滑向 $0$。从 $1.1$ 出发：$1.21$，$1.464$，
$2.144$，$4.595$——逃向无穷。\omterm{pb:g10:functions:1}{不动点} $0$
吸引，$1$ 排斥：起点相差毫厘，结局
却截然相反。

\textbf{11.} $7 \to 22 \to 11 \to 34 \to 17 \to 52 \to 26 \to
13 \to 40 \to 20 \to 10 \to 5 \to 16 \to 8 \to 4 \to 2 \to 1$:
共十六步，最高达 $52$。

\textbf{12.} $27 \to 82 \to 41 \to 124 \to 62 \to 31 \to 94 \to
47 \to 142 \to 71 \to 214 \to \dots$——走过十步后仍在攀升，
走向最高值 $9\,232$，并经历 $111$ 步的
“飞行”。由此可见，一条仅两行的规则，竟能产生
让人一眼难以预料的行为——相邻的数（$26$、$27$、$28$）
的飞行过程截然不同。

\textbf{13.} $2$ 的幂沿其整个阶梯逐次减半：
$2^k \to 2^{k-1} \to \dots \to 2 \to 1$。在最底端，
$h(1) = 4$，$h(4) = 2$，$h(2) = 1$：飞行便进入循环
$4 \to 2 \to 1 \to 4$。\omterm{pb:g10:functions:1}{不动点}需满足 $\frac n2 = n$
（仅 $n = 0$）或 $3n + 1 = n$（为负数）：正\omterm{def:g10:numbers:sets}{整数}中一个也没有
——观察到的飞行过程并不止于\omterm{pb:g10:functions:1}{不动点}
而是落入那个小小的\emph{循环}。

\textbf{14.} 反例要么是一条
无限增长（永不回到 $1$）的飞行过程，或是一个
与 $4 \to 2 \to 1$ 不相交的第二个循环。截至
$10^{20}$ 的验证仍留下无穷多未经验证的数——恰与
初中分册“火柴棒神谕”周末问题中圆形区域的教训相同：
在有限多个情形上相符，对全部情形什么也证明不了。

\textbf{15.} 对于海伦的机器，我们拥有\emph{结构}——一
个\omterm{def:g10:algebra:expand}{因式分解}的差（第 6 题）揭示了变化情况、
\omterm{def:g10:functions:extrema}{最小值}与夹逼——而冰雹数的奇偶转换规则
并不提供这样的抓手：无人找到能够
驯服它的结构。

\textbf{16.} $\sqrt{x - 3}$：需 $x - 3 \geq 0$：
$\intco{3}{+\infty}$。$\frac{1}{x^2 - 9}$：需
$x \neq \pm 3$：$\R$ 去掉 $-3$ 和 $3$。
$\sqrt{9 - x^2}$：需 $x^2 \leq 9$：$\intcc{-3}{3}$。

\textbf{17.} 由 $f(x) = 3$ 得 $x + \frac2x = 6$，于是
$x^2 - 6x + 2 = 0$，即 $(x - 3)^2 = 7$：
$x = 3 - \sqrt7$ 或 $x = 3 + \sqrt7$（两者均为正：共两个精确
\omterm{def:g10:functions:function}{原像}）。

\textbf{18.} 石子在空中时的实际意义：
$H(t) = 5t(4 - t) \geq 0$，对 $t \in \intcc{0}{4}$ 成立。配方得：
$H(t) = 20 - 5(t - 2)^2$：\omterm{def:g10:functions:extrema}{最大值} $20$ m，在 $t = 2$ s 时取得。
变化情况：在 $\intcc{0}{2}$ 上从 $0$ \omterm{def:g10:functions:variations}{递增}到 $20$，
在 $\intcc{2}{4}$ 上\omterm{def:g10:functions:variations}{递减}回 $0$。

\textbf{19.} $20t - 5t^2 \geq 15 \iff t^2 - 4t + 3 \leq 0 \iff
(t - 1)(t - 3) \leq 0 \iff t \in \intcc{1}{3}$：整整
两秒内高于 $15$ m。

\textbf{20.} 每台机器都有一个\omterm{def:g10:functions:function}{定义域}（规则
允许使用的范围）、一条规则（公式或奇偶切换），以及一个
能完整展示它的\omterm{def:g10:functions:graph}{图像}或表格。对每台机器，我们都提出了两个
大问题：\emph{重复导向何处}（\omterm{pb:g10:functions:1}{不动点}、
循环、吸引——对海伦的机器已解决，对冰雹数
尚无定论）；另一个是\emph{输出如何随输入变化}
（变化情况、极值——石子的最高点、海伦的\omterm{def:g10:functions:extrema}{最小值}）。
后续各章将深化这两个问题：接下来学习基本\omterm{def:g10:functions:function}{函数}，
而在以后的学习中，将借助导数刻画变化，
借助极限确认迭代的归宿。
\end{solution}
